\chapter{The Ontology of Vim}
\label{ch:ontology}

Before any grammar can be introduced, its terms have to be named. A sentence about deleting a word presupposes that both "deletion" and "word" already denote something in particular; a command like \texttt{daw} is meaningless until the reader knows what a text object is, just as much as it is meaningless until they know what an operator is. This chapter is that naming. It does not yet explain how these entities combine (that is the task of Part II) or what can be done with them (Part III onward). It only establishes what exists: the vocabulary the rest of the book will assume without re-explaining.

The entities below are not of the same kind. Some are persistent containers that outlive any single edit (buffers, registers). Some are transient views onto those containers (windows, tabs). Some are addresses into text (marks) or classes of region within it (text objects). Some are not things at all but states the editor itself can be in (modes). Keeping these categories distinct from the outset avoids a common source of confusion for newcomers, who often conflate a buffer with the window that happens to be displaying it, or a mark with the text object it can help construct.

\section{Buffers}

A buffer is Vim's in-memory representation of a file's content. When a file is opened, its contents are read into a buffer; edits are made to that buffer, not to the file on disk, until the buffer is explicitly written back with a command like \texttt{:w}. This distinction, buffer versus file, is not a technicality. A buffer can exist with no corresponding file at all (an empty scratch buffer created with \texttt{:enew}), and a single file can be associated with only one buffer at a time even if that buffer is currently displayed in several windows simultaneously. Every open file corresponds to exactly one buffer; every buffer does not necessarily correspond to a file.

Buffers persist independently of whether they are currently visible. A buffer can be "hidden," meaning it retains its content and any unsaved changes while no window is currently displaying it. This is what makes it possible to have dozens of files open in a single Vim session while only ever looking at one or two of them at a time: the buffer list is a superset of what is presently on screen. Chapter 12 covers the commands for listing, switching between, and managing buffers directly; for now, the essential fact is that a buffer is Vim's unit of "a file's content, currently loaded," and it is the thing every other entity in this chapter ultimately operates on or displays.

\section{Windows}

A window is a viewport onto a buffer. It occupies a rectangular region of the screen, has its own cursor position, its own scroll position, and its own local settings for certain options, but the buffer it displays is shared state: editing text through one window immediately changes what any other window displaying the same buffer shows, because there is only one buffer, not one per window. Vim allows the screen to be divided into multiple windows, arranged by horizontal or vertical splits (Chapter 22), which is what makes it possible to view two different regions of the same file, or two entirely different files, side by side.

The buffer/window distinction is the single most useful conceptual separation in this chapter for a newcomer to internalize early, because it explains behavior that otherwise looks surprising: closing a window with \texttt{:q} does not delete the buffer it was displaying, provided that buffer is still open in another window or has unsaved changes Vim is unwilling to discard; and opening the same file in two separate windows does not produce two independent copies of it to reconcile later, because both windows are looking at the same underlying buffer the entire time.

\section{Tabs}

A tab page, confusingly for anyone arriving from a browser's notion of "tabs," is not a container for a single file. It is a container for an arrangement of windows: a tab page can itself be split into several windows, each showing its own buffer. Switching tabs with \texttt{gt} or \texttt{gT} does not switch which file is open; it switches which saved window layout is currently visible. A user who wants to keep one project's files in one arrangement of splits and a second project's files in a completely different arrangement, and to switch between those two arrangements as a unit, is the intended use case for tabs. Chapter 23 covers this in more depth, including how tabs relate to sessions, which allow an entire multi-tab, multi-window, multi-buffer layout to be saved to disk and restored later.

\section{Files}

The relationship between a file (a named object on disk) and a buffer (Vim's loaded representation of that file's content) has already been introduced above, but it is worth stating the asymmetry explicitly, because several later chapters depend on it. A buffer can be modified without the underlying file changing at all, for as long as the user chooses not to write it. A file can change on disk, by some other process entirely, while a buffer holding a now-stale copy of its old content remains open and blissfully unaware, until Vim is prompted to check. Vim is, at its core, an editor of buffers that happens to load them from and save them to files; it is not, at the moment-to-moment level of an editing session, an editor of files directly. This distinction is also what makes constructs like scratch buffers, and Vim's ability to read a shell command's output directly into a buffer without ever touching disk (\texttt{:r !command}), coherent rather than exceptional: a buffer never actually required a file to exist in the first place.

\section{Registers}

A register is a small, named piece of persistent storage, holding a span of text, that survives across edits until it is next overwritten. Every deletion, change, and yank in Vim writes to some register, whether the user names one explicitly or not; the unnamed register (accessed as \texttt{"}) is what an unadorned \texttt{p} or \texttt{P} reads from, and it is silently updated by nearly every deleting or yanking command. Named registers, addressed with a letter (\texttt{"a} through \texttt{"z}), give the user a way to hold more than one piece of text at once, deliberately, rather than being limited to whatever the most recent operation happened to leave behind.

Registers are not merely a clipboard, and treating them as one obscures most of what makes them useful. There are numbered registers that automatically cycle through recent deletions, letting older deleted text be retrieved even after several more deletions have occurred; a read-only register holding the contents of the most recent search pattern; an expression register that evaluates arbitrary Vimscript and inserts the result; and, on systems where Vim is compiled with clipboard support, registers that read from and write to the operating system's own clipboard, bridging Vim's internal storage with the rest of the desktop environment. Chapter 16 treats every register in detail; the point to establish here is only that "register" denotes a general mechanism for named, persistent text storage, of which the familiar copy-paste clipboard is one instance rather than the whole story.

\section{Marks}

A mark is a saved position within a buffer: a specific line and column, given a one-character name, to which the cursor can later be returned directly. Local marks, named with a lowercase letter (\texttt{ma} sets mark \texttt{a}), are specific to the buffer in which they were set. Global marks, named with an uppercase letter, are specific to a position in a particular file but are addressable from any buffer, which makes them the mechanism for jumping directly to a remembered location in an entirely different file without first having to open and navigate to it manually.

Marks matter to this book for a reason beyond simple navigation: they are one of the things a motion can target, which means marks are among the building blocks operators compose with. The command \texttt{d\textquotesingle a} deletes from the cursor to mark \texttt{a}, treating the mark as a motion target in exactly the same grammatical slot a word or a search pattern would occupy. This is examined properly once the operator-motion grammar is introduced in Part II; here it is enough to know that a mark is a named, addressable position, and that Vim also maintains several marks automatically (the position before the last jump, the boundaries of the last changed or yanked text) without the user ever having to set them by hand, covered fully in the chapter on Marks in Part IV.

\section{Text Objects}

A text object denotes a structured region of text, defined not by explicit start and end positions the user specifies but by the syntactic or semantic unit it belongs to: a word, a quoted string, the contents of a pair of parentheses, a paragraph, a sentence. Where a motion moves the cursor to a position, a text object identifies a span, generally without moving the cursor itself when used as the argument to a command; \texttt{iw} does not mean "go to the word," it means "the inner word, wherever the cursor currently sits within it."

Text objects come in two forms for most of the structures they describe: an inner form (\texttt{i}) that selects the content of the structure without its delimiters, and an around form (\texttt{a}) that includes them. \texttt{ci"} changes the text inside a pair of double quotes without touching the quote marks themselves; \texttt{ca"} changes the quoted text and the quote marks together. This inner/around distinction recurs across nearly every text object Vim defines and is worth internalizing as a general pattern rather than memorizing separately for each one. Chapter 10 catalogs the full set; this section's purpose is only to place "text object" in the ontology as a distinct kind of thing from a motion, a mark, or a register, because later chapters will depend on the reader not conflating the three.

\section{Modes}

A mode is a state the editor itself is in, which determines what the same keystroke means. This was introduced conceptually in Chapter 1; here it is enough to enumerate what modes exist, without yet explaining their internal behavior in full, which is instead taken up as each mode becomes relevant to the surrounding material: Insert and Replace mode in the next part's chapter on Inserting and Replacing, Visual mode as it becomes relevant to later transformations, and Command-line mode throughout Part V. Vim distinguishes at minimum Normal mode (the default, for issuing commands), Insert mode (for typing text directly into the buffer), Visual mode (for selecting a region before acting on it, itself subdivided into character-wise, line-wise, and block-wise variants), Command-line mode (for entering Ex commands, searches, and filter expressions), and Replace mode (for overtyping existing characters rather than inserting new ones ahead of them). Terminal mode, present in Vim versions with an embedded terminal emulator, is a further mode specific to that feature. What unifies all of these, and what justifies treating "mode" as a first-class entity in this ontology rather than an implementation detail, is that the meaning of every subsequent keystroke in this book is conditioned on which mode is active at the moment it is pressed, and a great deal of Vim's apparent inconsistency to new users evaporates once this conditioning is made explicit rather than left implicit.

\section{Commands}

A command, in the broadest sense used in this book, is any complete instruction issued to Vim, whether it is a single Normal-mode keystroke, a composed sequence like \texttt{daw}, or a line typed in Command-line mode beginning with a colon. This is a deliberately broad category, and it is broad on purpose: one of the central claims of Part II is that Normal-mode key sequences and Ex commands are not two unrelated command languages coexisting inside one editor, but two surfaces of the same underlying grammar, addressed differently (one by direct keystroke, one by typed line) but sharing much of the same vocabulary of ranges, patterns, and targets. This chapter does not yet defend that claim; it only flags, as part of naming what exists, that "command" is being used here to cover both surfaces rather than being reserved for one of them.

\section{The Help System}

Vim ships with an extensive, cross-referenced help system, accessed with \texttt{:help} (or its abbreviation \texttt{:h}), covering every command, option, and feature described anywhere in this book, and a great many that are not. It is included here in the ontology chapter, rather than left as an afterthought, because it is not merely a manual bolted onto the editor; it is itself a Vim buffer, opened in a split window, navigable with the same motions, searches, and jumps as any other buffer, and containing hyperlink-like tags (rendered distinctly, and followed with \texttt{Ctrl-]}) that let the reader jump directly from a mention of a related command to that command's own help entry. A reader who has internalized nothing else from this chapter except how to move \texttt{:help} out of the way, follow a tag, and jump back with \texttt{Ctrl-o} already has a self-sufficient way to go beyond anything covered in this book, using the exact editing skills the rest of the book is teaching to navigate the documentation of the editor itself.

\section{Summary}

Nine entities have been named: buffers as the loaded content of a file, windows as viewports onto a buffer, tabs as saved arrangements of windows, files as the on-disk objects buffers are loaded from and written to, registers as named persistent text storage, marks as named positions, text objects as structured regions, modes as states that condition what a keystroke means, and commands as the general term for complete instructions in either the Normal-mode or Ex surface of Vim's grammar. None of these has yet been shown in combination with any other; that combination, and the rules that govern it, is the subject of Part II.
